\section{Related Work}
\label{related}
 Generative modeling for LiDAR has seen considerable development with the advent of several pattern recognition systems - diffusion models, generative models, etc. They use LiDAR point cloud or range image representation for underlying pattern recognition and modelling of LiDAR. Generating occluded static parts of scenes has been explored for several modalities: images, point clouds and recently LiDAR. 
 via synthetic-to-real transfer. 
 (\citet{wang2022coarse}) propose a coarse-to-fine network for dynamic-to-static image translation, without any segmentation information and show impressive results. 
 (\citet{ahn2023unsupervised}) propose to inpaint a dynamic semantic segmentation mask by replacing dynamic objects based labels with static object labels.
  For point clouds, (\citet{fei2023dctr}) address the problem of unstructuerd point prediction in lcoal patches and the inefficient utilization of local features. They propose a dual channel transformer that leverages point and channel-wise attention to generate locally compact and structured point clouds. (\citet{li2023progressive}) model the point cloud generation process a two-level hierarchy of distributions to learn the pattern of the point clouds. They propose a hierarchical consistency variational autoencoder that ensures that  the fine-grained parts and underlying shape and can be reconstructed with high precision. \citet{zhang2023learning} alleviate distortion of geometric details by formulating the  point cloud completion as a geometric transformation problem. They exploit repetitive geometric structures in 3D shapes to generate complete shapes.
  (\citet{zhang2021unsupervised}) uses GAN inversion for point cloud completion and shows remarkable results on standard shape datasets. (\citet{kumar2021dynamic}) achieves static translations for LiDAR scans using an adversarially trained model. They implicitly do not use segmentation information but have a variant that uses it for static translation. (\citet{yang2021semantic}) approach scene completion for LiDAR with the assistance of semantic-segmentation. (\citet{rist2021semantic}) attempt semantic scene completion using local Deep Implicit Functions as a learning based method. Their method focuses a continuous scene pattern representation unlike existing voxelization based works - this helps to encoder accurate semantic and geometric properties. (\citet{xia2023scpnet}) also perform semantic scene completion using a novel knowledge distillation objective - Dense-to-Sparse Knowledge Distillation.  

  
  For LiDAR, (\citet{zyrianov2022learning}) introduce LiDARGen that generates realistic high quality LiDAR samples using a stochastic denoising process that leverages the score-matching energy based model. They suffer from a major drawback - a high sampling rate of 20s/scan. This is infeasible for real world navigaiton that require a new scan within 100ms - freuqnecy of 10 Hz. (\citet{zhang2023nerf}) introduced NeRF-LiDAR, a NeRF based representation of scenes and use them to generate LiDAR scans. For a given LiDAR scene, NeRF-LiDAR requires multi-view images and semantic segmentation labels , both of which are unavailable in our setting. (\citet{xiong2023learning}) present UltraLiDAR that learns a discrete pattern of LiDAR. For LiDAR completion, the authors rely on explicit intervention to identify the codes for separate objects which have to be perturbed to obtain the resultant LiDAR.

Most static translation methods for point clouds discussed above focus on enhancing the power of the generator. Our focus is to enable the discriminator to learn precise discriminatory features between static and dynamic LIDAR scans - moving and movable objects.

\textbf{Contrastive Learning:}
Contrastive Learning focus on learning rich patterns in data by grouping samples with the same labels together while distancing out dissimilar ones.
Some well-known contrastive losses include Maximum Margin Contrastive Loss (\citet{hadsell2006dimensionality}), Triplet Loss (\citet{weinberger2009distance}) and multi-class N-pair loss (\citet{sohn2016improved}).
Triplet Loss operates on a triplet of latent vectors. Two vectors (the anchor and positive) share the same label while the third (negative) belongs to a different label. The objective of triplet loss is to find patterns in the dataset that keep the positive sample within an $\alpha$-neighborhood of the anchor while forcing the negative sample out of the $\alpha$-neighborhood, $\alpha$ being a margin hyperparamter.
The multi-class N-pair loss  is a generalization of the triplet loss.
Contrastive Learning with hard negatives (\citet{robinson2020contrastive}) involves selecting negative samples that are difficult to distinguish from the anchor sample. Several approaches (\citet{xie2022delving}) deploy hard negative mining for contrastive learning in order to learn better representations for downstream tasks.